% !TeX program = xelatex
\documentclass[12pt,a4paper]{ctexart}

\usepackage{amsmath,amssymb,mathtools}
\usepackage{geometry}
\usepackage{enumitem}

\geometry{
  left=2.5cm,
  right=2.5cm,
  top=2.5cm,
  bottom=2.5cm
}

\setlength{\parindent}{2em}
\setlength{\parskip}{0.3em}
\allowdisplaybreaks

\newcommand{\dd}{\,\mathrm d}
\newcommand{\R}{\mathbb R}

\title{\textbf{第二次作业（第二周 9.24）}\\[0.5em]
\large 题目与解答}
\author{}
\date{}

\begin{document}

\maketitle

\noindent
\textbf{说明：}
原题第 1、2 题中
\[
P\in\mathbb R^2
\]
应为
\[
P\subset\mathbb R^2,
\]
以下已作相应修正。

% ============================================================
\section*{1.\ 含参变量反常积分的基本性质}

\subsection*{题目}

设
\[
P=[a,\omega)\times Y\subset\R^2,
\]
函数
\[
f(x,y)\in C(P),
\]
函数 \(g(x)\) 是 \([a,\omega)\) 上的局部可积函数，即对任意
\[
[a,b]\subset[a,\omega),
\]
都有
\[
g\in R[a,b].
\]
考虑含参变量反常积分
\[
F(y)=\int_a^\omega f(x,y)g(x)\dd x.
\]

证明：

\begin{enumerate}[label=(\arabic*),leftmargin=3em]
\item
设 \(Y=[c,d]\)。若 \(F(y)\) 在 \(Y\) 上一致收敛，则
\[
F\in C[c,d].
\]

\item
设 \(Y=[c,d]\)。若
\[
\partial_y f(x,y)\in C(P),
\]
且
\[
\int_a^\omega \partial_y f(x,y)g(x)\dd x
\]
在 \(Y\) 上一致收敛，而 \(F(y)\) 在某一点
\[
y_0\in Y
\]
收敛，则 \(F(y)\) 在 \(Y\) 上收敛且一致收敛，并且
\[
F\in C^1[c,d],
\]
且
\[
F'(y)
=
\int_a^\omega
\partial_y f(x,y)g(x)\dd x.
\]

\item
设 \(Y=[c,d]\)。若 \(F(y)\) 在 \(Y\) 上一致收敛，则
\[
F\in R[c,d],
\]
且
\[
\int_c^d F(y)\dd y
=
\int_a^\omega
\left(
\int_c^d f(x,y)g(x)\dd y
\right)\dd x.
\]

\item
设
\[
Y=[c,\widetilde\omega).
\]
若定义 \(F(y)\) 的反常积分在任意
\[
[c,d]\subset[c,\widetilde\omega)
\]
上一致收敛，且
\[
\Phi(x)
=
\int_c^{\widetilde\omega}
f(x,y)g(x)\dd y
\]
在任意
\[
[a,b]\subset[a,\omega)
\]
上一致收敛，并且下列两个绝对值累次积分至少有一个存在：
\[
\int_c^{\widetilde\omega}
\dd y
\int_a^\omega
|f(x,y)g(x)|\dd x,
\]
\[
\int_a^\omega
\dd x
\int_c^{\widetilde\omega}
|f(x,y)g(x)|\dd y,
\]
证明
\[
\int_c^{\widetilde\omega}F(y)\dd y
=
\int_a^\omega
\dd x
\int_c^{\widetilde\omega}
f(x,y)g(x)\dd y.
\]
\end{enumerate}

\subsection*{解答}

对 \(b<\omega\)，记
\[
F_b(y)
=
\int_a^b f(x,y)g(x)\dd x.
\]

在任意紧矩形 \([a,b]\times[c,d]\) 上，由第一次作业第 1 题
关于含参变量常义积分连续性的结论，
\[
F_b\in C[c,d].
\]

\medskip

\noindent
\textbf{(1)}

由假设，反常积分
\[
F(y)=\int_a^\omega f(x,y)g(x)\dd x
\]
在 \(Y=[c,d]\) 上一致收敛，也就是说
\[
F_b(y)\longrightarrow F(y)
\qquad (b\uparrow\omega)
\]
在 \([c,d]\) 上一致成立。

每个 \(F_b\) 都是连续函数，而连续函数列的一致极限仍连续，所以
\[
\boxed{F\in C[c,d].}
\]

\medskip

\noindent
\textbf{(2)}

记
\[
G_b(y)
=
\int_a^b
\partial_y f(x,y)g(x)\dd x.
\]
由于
\[
\partial_y f\in C(P),
\]
由第一次作业第 1 题关于积分号下求导的结论，
\[
F_b'(y)=G_b(y).
\]

由题设，
\[
G_b(y)
\longrightarrow
G(y)
:=
\int_a^\omega
\partial_y f(x,y)g(x)\dd x
\]
在 \(Y=[c,d]\) 上一致收敛。

又因为 \(F(y_0)\) 收敛，所以存在常数 \(L\)，使得
\[
F_b(y_0)\longrightarrow L
\qquad (b\uparrow\omega).
\]

对任意 \(y\in[c,d]\)，由微积分基本定理，
\[
F_b(y)-F_b(y_0)
=
\int_{y_0}^yG_b(t)\dd t.
\]
因此
\[
F_b(y)
=
F_b(y_0)
+
\int_{y_0}^yG_b(t)\dd t.
\]

令 \(b\uparrow\omega\)，得到候选极限
\[
F_*(y)
=
L+\int_{y_0}^yG(t)\dd t.
\]
并且
\begin{align*}
|F_b(y)-F_*(y)|
&\le
|F_b(y_0)-L|
+
\left|
\int_{y_0}^y
\bigl(G_b(t)-G(t)\bigr)\dd t
\right|\\
&\le
|F_b(y_0)-L|
+
(d-c)
\sup_{t\in[c,d]}|G_b(t)-G(t)|.
\end{align*}
右端与 \(y\) 无关并趋于 \(0\)，故
\[
F_b\longrightarrow F_*
\]
在 \([c,d]\) 上一致收敛。

这说明对每个 \(y\in[c,d]\)，原反常积分 \(F(y)\) 均收敛，
且
\[
F(y)=F_*(y).
\]
所以 \(F(y)\) 在 \(Y\) 上一致收敛。

另一方面，每个 \(G_b\) 连续，而 \(G_b\to G\) 一致，
故
\[
G\in C[c,d].
\]
由
\[
F(y)
=
L+\int_{y_0}^yG(t)\dd t
\]
可知
\[
F\in C^1[c,d]
\]
且
\[
F'(y)=G(y).
\]
因此
\[
\boxed{
F'(y)
=
\int_a^\omega
\partial_y f(x,y)g(x)\dd x.
}
\]

\medskip

\noindent
\textbf{(3)}

由第 (1) 问，
\[
F\in C[c,d],
\]
因而当然
\[
F\in R[c,d].
\]

对于任意 \(b<\omega\)，在有限矩形
\[
[a,b]\times[c,d]
\]
上由第一次作业第 1 题的换序结论，于是
\begin{align*}
\int_c^dF_b(y)\dd y
&=
\int_c^d
\left(
\int_a^b f(x,y)g(x)\dd x
\right)\dd y\\
&=
\int_a^b
\left(
\int_c^d f(x,y)g(x)\dd y
\right)\dd x.
\end{align*}

由于
\[
F_b\longrightarrow F
\]
在 \([c,d]\) 上一致收敛，
所以可以在有限区间上逐项积分：
\[
\lim_{b\uparrow\omega}
\int_c^dF_b(y)\dd y
=
\int_c^dF(y)\dd y.
\]
故
\[
\lim_{b\uparrow\omega}
\int_a^b
\left(
\int_c^d f(x,y)g(x)\dd y
\right)\dd x
=
\int_c^dF(y)\dd y.
\]
这同时说明右边关于 \(x\) 的反常积分存在，并且
\[
\boxed{
\int_c^dF(y)\dd y
=
\int_a^\omega
\left(
\int_c^d f(x,y)g(x)\dd y
\right)\dd x.
}
\]

\medskip

\noindent
\textbf{(4)}

记
\[
H(x,y)=f(x,y)g(x).
\]
对 \(b<\omega\)、\(d<\widetilde\omega\)，令
\[
Q(b,d)
=
\int_a^b\dd x\int_c^dH(x,y)\dd y,
\]
以及
\[
M(b,d)
=
\int_a^b\dd x\int_c^d|H(x,y)|\dd y.
\]
在有限矩形上由第一次作业第 1 题的换序结论也可写成
\[
Q(b,d)
=
\int_c^d\dd y\int_a^bH(x,y)\dd x,
\]
\[
M(b,d)
=
\int_c^d\dd y\int_a^b|H(x,y)|\dd x.
\]

根据题设，两个绝对值累次积分中至少有一个存在。
无论是哪一个存在，都有某个常数 \(C<\infty\)，使得对所有
\(b<\omega,d<\widetilde\omega\)，
\[
0\le M(b,d)\le C.
\]

定义
\[
M_b
=
\lim_{d\uparrow\widetilde\omega}M(b,d),
\]
以及
\[
N_d
=
\lim_{b\uparrow\omega}M(b,d).
\]
由于 \(M(b,d)\) 对 \(b,d\) 都单调不减且有统一上界，
这些极限都存在。

再令
\[
M_*=\sup_{b,d}M(b,d).
\]
则
\[
M_b\uparrow M_*
\qquad (b\uparrow\omega),
\]
且
\[
N_d\uparrow M_*
\qquad (d\uparrow\widetilde\omega).
\]

先考察关于 \(x\) 的累次积分。
由于
\[
\int_c^{\widetilde\omega}H(x,y)\dd y=\Phi(x)
\]
在每个紧区间 \([a,b]\) 上一致收敛，所以
\[
\int_a^b\Phi(x)\dd x
=
\lim_{d\uparrow\widetilde\omega}Q(b,d).
\]
若 \(b_2>b_1\)，则
\begin{align*}
\left|
\int_{b_1}^{b_2}\Phi(x)\dd x
\right|
&=
\left|
\lim_{d\uparrow\widetilde\omega}
\bigl(Q(b_2,d)-Q(b_1,d)\bigr)
\right|\\
&\le
\lim_{d\uparrow\widetilde\omega}
\bigl(M(b_2,d)-M(b_1,d)\bigr)\\
&=
M_{b_2}-M_{b_1}\\
&\le
M_*-M_{b_1}.
\end{align*}
当 \(b_1\uparrow\omega\) 时，右端趋于 \(0\)。
由反常积分的 Cauchy 判别法，
\[
\int_a^\omega\Phi(x)\dd x
\]
存在。

完全类似地，由 \(F(y)\) 在每个
\([c,d]\) 上的一致收敛可得
\[
\int_c^dF(y)\dd y
=
\lim_{b\uparrow\omega}Q(b,d).
\]
若 \(d_2>d_1\)，则
\[
\left|
\int_{d_1}^{d_2}F(y)\dd y
\right|
\le
N_{d_2}-N_{d_1}
\le
M_*-N_{d_1}.
\]
因此
\[
\int_c^{\widetilde\omega}F(y)\dd y
\]
也存在。

最后说明两者相等。

若 \(b_2\ge b_1\)、\(d_2\ge d_1\)，则
\begin{align*}
|Q(b_2,d_2)-Q(b_1,d_1)|
&\le
|Q(b_2,d_2)-Q(b_1,d_2)|\\
&\qquad
+
|Q(b_1,d_2)-Q(b_1,d_1)|\\
&\le
M_{b_2}-M_{b_1}
+
N_{d_2}-N_{d_1}\\
&\le
(M_*-M_{b_1})
+
(M_*-N_{d_1}).
\end{align*}
当
\[
b_1\uparrow\omega,\qquad
d_1\uparrow\widetilde\omega
\]
时，右端趋于 \(0\)。

因此 \(Q(b,d)\) 在
\[
b\uparrow\omega,\qquad
d\uparrow\widetilde\omega
\]
时存在共同极限，故两个累次极限相等：
\begin{align*}
\int_c^{\widetilde\omega}F(y)\dd y
&=
\lim_{d\uparrow\widetilde\omega}
\lim_{b\uparrow\omega}Q(b,d)\\
&=
\lim_{b\uparrow\omega}
\lim_{d\uparrow\widetilde\omega}Q(b,d)\\
&=
\int_a^\omega\Phi(x)\dd x.
\end{align*}
即
\[
\boxed{
\int_c^{\widetilde\omega}F(y)\dd y
=
\int_a^\omega
\dd x
\int_c^{\widetilde\omega}
f(x,y)g(x)\dd y.
}
\]

% ============================================================
\section*{2.\ 非负函数的累次反常积分}

\subsection*{题目}

设
\[
P=[a,\omega)\times[c,\widetilde\omega)
\subset\R^2,
\]
非负函数
\[
f(x,y)\in C(P).
\]
考虑含参变量反常积分
\[
F(y)=\int_a^\omega f(x,y)\dd x,
\qquad
\Phi(x)=\int_c^{\widetilde\omega}f(x,y)\dd y.
\]

若
\[
F\in C[c,\widetilde\omega),
\qquad
\Phi\in C[a,\omega),
\]
且下列两个累次积分至少有一个存在：
\[
\int_c^{\widetilde\omega}
\dd y
\int_a^\omega f(x,y)\dd x,
\]
\[
\int_a^\omega
\dd x
\int_c^{\widetilde\omega}f(x,y)\dd y,
\]
证明另一个也存在，且两者相等。

\subsection*{解答}

记
\[
F_b(y)=\int_a^b f(x,y)\dd x,
\qquad
\Phi_d(x)=\int_c^d f(x,y)\dd y.
\]

固定任意紧区间 \([c,d]\subset[c,\widetilde\omega)\)。
将第一次作业第 5 题的结论用于非负函数 \(f\)，由
\(F\in C[c,d]\) 可得
\[
F_b\longrightarrow F
\]
在 \([c,d]\) 上一致收敛。交换 \(x,y\) 的角色，同理，对任意
紧区间 \([a,b]\subset[a,\omega)\)，由 \(\Phi\in C[a,b]\) 可得
\[
\Phi_d\longrightarrow\Phi
\]
在 \([a,b]\) 上一致收敛。

不妨假设
\[
A
:=
\int_c^{\widetilde\omega}F(y)\dd y
=
\int_c^{\widetilde\omega}
\dd y
\int_a^\omega f(x,y)\dd x
<\infty.
\]
我们证明另一个累次积分存在且等于 \(A\)。

对任意 \(a<b<\omega\)，由 \(\Phi_d\to\Phi\) 在
\([a,b]\) 上一致收敛以及第一次作业第 1 题的换序结论，
\begin{align*}
\int_a^b\Phi(x)\dd x
&=
\lim_{d\uparrow\widetilde\omega}
\int_a^b\Phi_d(x)\dd x\\
&=
\lim_{d\uparrow\widetilde\omega}
\int_a^b\dd x\int_c^d f(x,y)\dd y\\
&=
\lim_{d\uparrow\widetilde\omega}
\int_c^d\dd y\int_a^b f(x,y)\dd x.
\end{align*}
因为 \(f\ge0\)，有
\[
\int_a^b f(x,y)\dd x
\le
F(y),
\]
故
\[
\int_a^b\Phi(x)\dd x
\le A.
\]

另一方面，\(\Phi(x)\ge0\)，所以
\[
\int_a^b\Phi(x)\dd x
\]
随 \(b\uparrow\omega\) 单调不减且有上界 \(A\)。
因此极限存在，即
\[
B
:=
\int_a^\omega\Phi(x)\dd x
\]
存在，并且
\[
B\le A.
\]

下面证明反向不等式。

对任意 \(d<\widetilde\omega\)，由
\(F_b\to F\) 在 \([c,d]\) 上一致收敛，
\begin{align*}
\int_c^dF(y)\dd y
&=
\lim_{b\uparrow\omega}
\int_c^dF_b(y)\dd y\\
&=
\lim_{b\uparrow\omega}
\int_c^d\dd y\int_a^b f(x,y)\dd x\\
&=
\lim_{b\uparrow\omega}
\int_a^b\dd x\int_c^d f(x,y)\dd y.
\end{align*}
又因 \(f\ge0\)，
\[
\int_c^d f(x,y)\dd y
\le
\Phi(x),
\]
所以
\[
\int_c^dF(y)\dd y
\le
B.
\]
令 \(d\uparrow\widetilde\omega\)，得到
\[
A\le B.
\]
结合 \(B\le A\)，可得
\[
A=B.
\]

因此
\[
\boxed{
\int_c^{\widetilde\omega}
\dd y
\int_a^\omega f(x,y)\dd x
=
\int_a^\omega
\dd x
\int_c^{\widetilde\omega}f(x,y)\dd y.
}
\]

若一开始假设右边的累次积分存在，只需交换 \(x,y\) 的角色，
同理可得相同结论。

% ============================================================
\section*{3.\ 两个振荡反常积分}

\subsection*{题目}

验证
\[
\int_0^\infty\frac{\cos x}{x^\alpha}\dd x
=
\frac{\pi}
{2\Gamma(\alpha)\cos\left(\dfrac{\pi\alpha}{2}\right)},
\qquad
\alpha\in(0,1),
\]
以及
\[
\int_0^\infty\frac{\sin x}{x^\beta}\dd x
=
\frac{\pi}
{2\Gamma(\beta)\sin\left(\dfrac{\pi\beta}{2}\right)},
\qquad
\beta\in(0,2).
\]

请给出详细过程并说明理由。

\subsection*{解答}

首先说明两个积分的收敛性。

对于
\[
\int_0^\infty\frac{\cos x}{x^\alpha}\dd x,
\qquad 0<\alpha<1,
\]
在 \(x=0\) 附近，
\[
\left|\frac{\cos x}{x^\alpha}\right|
\le \frac1{x^\alpha},
\]
而
\[
\int_0^1x^{-\alpha}\dd x<\infty.
\]
在无穷远处，\(x^{-\alpha}\) 单调趋于 \(0\)，
故由 Dirichlet 判别法，
\[
\int_1^\infty x^{-\alpha}\cos x\dd x
\]
收敛。

类似地，在 \(x=0\) 附近，
\[
\frac{\sin x}{x^\beta}
\sim x^{1-\beta},
\]
所以当且仅当 \(\beta<2\) 时局部可积。
在无穷远处，\(x^{-\beta}\downarrow0\)，
故当 \(\beta>0\) 时由 Dirichlet 判别法收敛。
因此题中给出的参数范围正是
\[
0<\alpha<1,
\qquad
0<\beta<2.
\]

\medskip

\noindent
\textbf{第一部分：计算余弦积分。}

令
\[
s=1-\alpha\in(0,1).
\]
对 \(\varepsilon>0\)，由 Gamma 函数的 Laplace 积分公式
\[
\int_0^\infty x^{s-1}e^{-zx}\dd x
=
\Gamma(s)z^{-s},
\qquad
\Re z>0,
\]
取
\[
z=\varepsilon-i,
\]
得到
\[
\int_0^\infty
x^{-\alpha}e^{-\varepsilon x}e^{ix}\dd x
=
\Gamma(1-\alpha)
(\varepsilon-i)^{-(1-\alpha)}.
\]
取实部：
\[
\int_0^\infty
e^{-\varepsilon x}
\frac{\cos x}{x^\alpha}\dd x
=
\Gamma(1-\alpha)
\Re\left(
(\varepsilon-i)^{-(1-\alpha)}
\right).
\]

由于原反常积分已经收敛，Abel 极限定理给出
\[
\lim_{\varepsilon\downarrow0}
\int_0^\infty
e^{-\varepsilon x}
\frac{\cos x}{x^\alpha}\dd x
=
\int_0^\infty
\frac{\cos x}{x^\alpha}\dd x.
\]

当 \(\varepsilon\downarrow0\) 时，
\[
\varepsilon-i
\longrightarrow
e^{-i\pi/2},
\]
于是
\begin{align*}
\int_0^\infty
\frac{\cos x}{x^\alpha}\dd x
&=
\Gamma(1-\alpha)
\cos\left(
\frac{(1-\alpha)\pi}{2}
\right)\\
&=
\Gamma(1-\alpha)
\sin\left(
\frac{\pi\alpha}{2}
\right).
\end{align*}

利用 Euler 反射公式
\[
\Gamma(\alpha)\Gamma(1-\alpha)
=
\frac{\pi}{\sin\pi\alpha},
\]
以及
\[
\sin\pi\alpha
=
2\sin\frac{\pi\alpha}{2}
\cos\frac{\pi\alpha}{2},
\]
可得
\begin{align*}
\Gamma(1-\alpha)
\sin\frac{\pi\alpha}{2}
&=
\frac{\pi}
{\Gamma(\alpha)\sin\pi\alpha}
\sin\frac{\pi\alpha}{2}\\
&=
\frac{\pi}
{2\Gamma(\alpha)
\cos\left(\dfrac{\pi\alpha}{2}\right)}.
\end{align*}
因此
\[
\boxed{
\int_0^\infty\frac{\cos x}{x^\alpha}\dd x
=
\frac{\pi}
{2\Gamma(\alpha)
\cos\left(\dfrac{\pi\alpha}{2}\right)},
\qquad 0<\alpha<1.
}
\]

\medskip

\noindent
\textbf{第二部分：计算正弦积分。}

先设
\[
0<\beta<1.
\]
同样由 Gamma 函数公式，
\[
\int_0^\infty
x^{-\beta}
e^{-\varepsilon x}e^{ix}\dd x
=
\Gamma(1-\beta)
(\varepsilon-i)^{-(1-\beta)}.
\]
取虚部并令 \(\varepsilon\downarrow0\)，由 Abel 极限定理得到
\begin{align*}
\int_0^\infty
\frac{\sin x}{x^\beta}\dd x
&=
\Gamma(1-\beta)
\sin\left(
\frac{(1-\beta)\pi}{2}
\right)\\
&=
\Gamma(1-\beta)
\cos\left(
\frac{\pi\beta}{2}
\right).
\end{align*}
利用反射公式，
\begin{align*}
\Gamma(1-\beta)
\cos\frac{\pi\beta}{2}
&=
\frac{\pi}
{\Gamma(\beta)\sin\pi\beta}
\cos\frac{\pi\beta}{2}\\
&=
\frac{\pi}
{2\Gamma(\beta)
\sin\left(\dfrac{\pi\beta}{2}\right)}.
\end{align*}
因此当 \(0<\beta<1\) 时结论成立。

下面考虑
\[
\beta=1.
\]
对 \(\varepsilon>0\)，令
\[
K_\varepsilon(a)
=
\int_0^\infty
e^{-\varepsilon x}
\frac{\sin(ax)}x\dd x.
\]
则
\[
K_\varepsilon'(a)
=
\int_0^\infty
e^{-\varepsilon x}\cos(ax)\dd x
=
\frac{\varepsilon}{\varepsilon^2+a^2}.
\]
又 \(K_\varepsilon(0)=0\)，所以
\[
K_\varepsilon(1)
=
\int_0^1
\frac{\varepsilon}{\varepsilon^2+a^2}\dd a
=
\arctan\frac1\varepsilon.
\]
令 \(\varepsilon\downarrow0\)，得到
\[
\int_0^\infty\frac{\sin x}{x}\dd x
=
\frac{\pi}{2},
\]
恰好与公式
\[
\frac{\pi}
{2\Gamma(1)\sin(\pi/2)}
=
\frac{\pi}{2}
\]
一致。

最后设
\[
1<\beta<2.
\]
对反常积分分部积分：
\[
u=\sin x,
\qquad
\dd v=x^{-\beta}\dd x,
\]
则
\[
v=\frac{x^{1-\beta}}{1-\beta}.
\]
由于
\[
x^{1-\beta}\sin x\longrightarrow0
\qquad (x\to0^+),
\]
因为此时
\[
x^{1-\beta}\sin x\sim x^{2-\beta}\to0,
\]
而且
\[
x^{1-\beta}\sin x\longrightarrow0
\qquad (x\to\infty),
\]
故边界项为零。于是
\begin{align*}
\int_0^\infty
\frac{\sin x}{x^\beta}\dd x
&=
\frac1{\beta-1}
\int_0^\infty
x^{1-\beta}\cos x\dd x\\
&=
\frac1{\beta-1}
\int_0^\infty
\frac{\cos x}{x^{\beta-1}}\dd x.
\end{align*}

因为
\[
0<\beta-1<1,
\]
可使用第一部分的结论：
\[
\int_0^\infty
\frac{\cos x}{x^{\beta-1}}\dd x
=
\frac{\pi}
{2\Gamma(\beta-1)
\cos\left(\dfrac{\pi(\beta-1)}2\right)}.
\]
注意
\[
\Gamma(\beta)
=
(\beta-1)\Gamma(\beta-1)
\]
以及
\[
\cos\left(
\frac{\pi(\beta-1)}2
\right)
=
\sin\left(
\frac{\pi\beta}{2}
\right),
\]
故
\[
\boxed{
\int_0^\infty
\frac{\sin x}{x^\beta}\dd x
=
\frac{\pi}
{2\Gamma(\beta)
\sin\left(\dfrac{\pi\beta}{2}\right)},
\qquad 0<\beta<2.
}
\]

\medskip

这里所用的 Abel 极限定理是：
若反常积分
\[
\int_0^\infty h(x)\dd x
\]
收敛，则
\[
\lim_{\varepsilon\downarrow0}
\int_0^\infty e^{-\varepsilon x}h(x)\dd x
=
\int_0^\infty h(x)\dd x.
\]
因此上述由指数衰减积分向原振荡反常积分取极限的步骤是合法的。

% ============================================================
\section*{4.\ Riemann--Lebesgue 引理的一个特殊形式}

\subsection*{题目}

设
\[
f\in R[a,b].
\]
证明
\[
\lim_{n\to\infty}
\int_a^b f(x)\sin nx\dd x=0.
\]

提示：
\[
\int_{\frac{2k\pi}{n}}^{\frac{(2k+2)\pi}{n}}
f(x)\sin nx\dd x
=
\int_{\frac{2k\pi}{n}}^{\frac{(2k+2)\pi}{n}}
\left[
f(x)-f\left(\frac{2k\pi}{n}\right)
\right]
\sin nx\dd x.
\]

\subsection*{解答}

因为 \(f\in R[a,b]\)，所以 \(f\) 在 \([a,b]\) 上有界。
设
\[
|f(x)|\le M.
\]

考虑长度为一个正弦周期的区间
\[
I_{k,n}
=
\left[
\frac{2k\pi}{n},
\frac{2(k+1)\pi}{n}
\right].
\]
由于
\[
\int_{I_{k,n}}\sin nx\dd x=0,
\]
故只要 \(I_{k,n}\subset[a,b]\)，就有
\begin{align*}
\int_{I_{k,n}}f(x)\sin nx\dd x
&=
\int_{I_{k,n}}
\left[
f(x)-f\left(\frac{2k\pi}{n}\right)
\right]
\sin nx\dd x.
\end{align*}
这正是题目所给的提示。

记 \(f\) 在 \(I_{k,n}\) 上的振幅为
\[
\operatorname{osc}(f;I_{k,n})
=
\sup_{x\in I_{k,n}}f(x)
-
\inf_{x\in I_{k,n}}f(x).
\]
则
\begin{align*}
\left|
\int_{I_{k,n}}f(x)\sin nx\dd x
\right|
&\le
\int_{I_{k,n}}
\left|
f(x)-f\left(\frac{2k\pi}{n}\right)
\right|\dd x\\
&\le
\operatorname{osc}(f;I_{k,n})
\,|I_{k,n}|\\
&=
\operatorname{osc}(f;I_{k,n})
\frac{2\pi}{n}.
\end{align*}

现在用所有位于 \([a,b]\) 内的完整周期区间
\(I_{k,n}\)，再加上左右两端不足一个周期的区间，
组成 \([a,b]\) 的一个分割 \(P_n\)。

该分割的网格满足
\[
\|P_n\|
\le
\frac{2\pi}{n}
\longrightarrow0.
\]

由于 \(f\) Riemann 可积，对任意网格趋于 \(0\) 的分割序列，
都有
\[
U(P_n,f)-L(P_n,f)\longrightarrow0.
\]
这里 \(U,L\) 分别表示 Darboux 上和与下和。

所有完整周期区间上的积分之和满足
\begin{align*}
\left|
\sum_{I_{k,n}\subset[a,b]}
\int_{I_{k,n}}
f(x)\sin nx\dd x
\right|
&\le
\sum_{I_{k,n}\subset[a,b]}
\operatorname{osc}(f;I_{k,n})|I_{k,n}|\\
&\le
U(P_n,f)-L(P_n,f).
\end{align*}

另一方面，左右两端至多各剩下一段长度小于
\(2\pi/n\) 的区间，因此这两个端部区间总长度小于
\[
\frac{4\pi}{n}.
\]
于是端部积分的绝对值不超过
\[
\frac{4\pi M}{n}.
\]

综上，
\[
\left|
\int_a^bf(x)\sin nx\dd x
\right|
\le
U(P_n,f)-L(P_n,f)
+
\frac{4\pi M}{n}.
\]
令 \(n\to\infty\)，右端趋于 \(0\)，故
\[
\boxed{
\lim_{n\to\infty}
\int_a^bf(x)\sin nx\dd x
=
0.
}
\]

这就是 Riemann--Lebesgue 引理在 Riemann 可积函数情形下的一个特殊形式。

\end{document}
